\section{Prefill 阶段}

\subsection{什么是 Prefill？}

\begin{importantbox}{Prefill = 处理输入 Prompt}
用户发送 prompt 后，模型需要一次性处理所有输入 token。这个阶段叫做 \textbf{Prefill}：
\begin{itemize}
  \item $Q$ 的长度 = prompt 长度 $N$（通常几百到几千）
  \item $K, V$ 的长度也是 $N$（因为是自注意力）
  \item 是一个大规模的 \textbf{矩阵-矩阵乘法}（GEMM）
  \item \textbf{计算密集型}（Compute-bound）
\end{itemize}
\end{importantbox}

\subsection{计算特性}

Prefill 阶段的核心计算是 $Q \times K^T$，这是一个 $N \times d_k$ 乘以 $d_k \times N$ 的矩阵乘法，结果为 $N \times N$。

\begin{itemize}
  \item 计算量：$O(N^2 \cdot d_k)$
  \item 数据量：$O(N \cdot d_k)$
  \item 计算/内存比（Arithmetic Intensity）很高 $\Rightarrow$ \textbf{计算密集型}
  \item GPU 的算力是瓶颈，内存带宽不是
\end{itemize}

\subsection{本章小结}

Prefill 是大矩阵乘法（GEMM），计算密集型，瓶颈在算力而非内存带宽。

